
\label{sec:alpha-dos-vias-t1-rev8}

The closure of \(\alpha\) does not begin with a metrological constant or with a
sought decimal string. It begins in the same enriched terminal state that
preserves the genealogy of the three channels. This section presents, within a single
causal chain, two routes to the scale of \(\alpha\) and a long prolongation
of the carry of the first. The remote prolongation is denoted by
\(\mathsf T_{\alpha,\infty}\); it constitutes neither a second route nor a third
derivation.

\subsection{Common terminal state and 2 coordinated internal readings}

Let
\[
 x_{\rm term}\in\mathcal X_{\rm TPK}^{\rm term,enr}
 \subseteq\mathcal X_{\rm TPK}^{[108]}.
\]
The terminal selector and the signed chart are two evaluations of the same state:
\begin{equation}
 x_{\rm term}
 \xrightarrow{(\pi_{\rm term},R_{\rm sgn})}
 \bigl(B_{\rm term},U_{12}^{\rm sgn}\bigr),
 \label{eq:alpha-rev8-doble-lectura-terminal}
\end{equation}
with
\[
 B_{\rm term}=
 \begin{pmatrix}
 021020\\001220\\100210
 \end{pmatrix}
\]
and
\[
 U_{12}^{\rm sgn}
 =(2378,1406,2479,-452,998,-551,-668,-204,-371,-322,-28,-997).
\]
The first reading retains the three selected terminal blocks; the
second retains the signed integer coordinate required for dodecaphase
closure. No factorization from
\(B_{\rm term}\) to \(U_{12}^{\rm sgn}\) exists or is needed. Drawing one would erase precisely the
fields of orientation, carry, boundary, and ledger that distinguish the two
readings. From the second, however, the reversible charts have been constructed
\[
 U_{12}^{\rm sgn}
 \xleftrightarrow{\ \mathcal H_{12}\ }K
 \xleftrightarrow{\ \mathcal E_{\rm dod}\ }
 (D_3K,D_4K,Q),
 \qquad
 \mathcal E_{\rm dod}=(D_3,D_4,Q).
\]
The domain of this closure is therefore the enriched TPK profile; a
raw trace of 108 positions is a forgetful projection and does not replace that
domain.

\subsection{Path A: dodecaphasic integer closure}

The Archimedean evaluations of the three channels and the canonical seal are
the twelve-block vectors
\[
\begin{aligned}
P={}&(141,592,653,589,793,238,462,643,383,279,502,884),\\
E={}&(718,281,828,459,045,235,360,287,471,352,662,497),\\
\Phi={}&(618,033,988,749,894,848,204,586,834,365,638,117),\\
K={}&(234,543,140,729,659,824,621,058,914,794,146,601).
\end{aligned}
\]
With the declared right boundary \(c_{13}=0\), Euclidean division
propagates from right to left:
\begin{equation}
\begin{aligned}
 s_m&=P_m+E_m-\Phi_m-K_m+c_{m+1},\\
 A_m&=s_m\bmod 1000,\\
 c_m&=\frac{s_m-A_m}{1000},
 \qquad m=12,11,\ldots,1.
\end{aligned}
\label{eq:alpha-rev8-via-entera}
\end{equation}
The output is
\[
 A=(007,297,352,569,283,800,997,285,105,472,380,663)
\]
and fixes the rational
\begin{equation}
 \alpha_{12}
 =0.007297352569283800997285105472380663.
 \label{eq:alpha-rev8-salida-corta}
\end{equation}
The abbreviated equality \(P+E-\Phi-K=A\) must be read in the chart of
integer concatenation. If
\[
 \iota(v_1,\ldots,v_{12})
 =\sum_{m=1}^{12}v_m1000^{12-m},
\]
then
\begin{equation}
 \boxed{\iota(P)+\iota(E)-\iota(\Phi)-\iota(K)=\iota(A)}.
 \label{eq:alpha-rev8-identidad-concatenada}
\end{equation}
In coordinates, the cocycle cannot be omitted:
\[
 P+E-\Phi-K+C^+-1000C=A,
 \qquad C^+=(c_2,\ldots,c_{13}).
\]
The termination \(\ldots|380|663\) belongs to this finite window and its
condition \(c_{13}=0\). The complete theorem, the table of the twelve carries, and
the inversion that recovers \(K\) reside in the preceding chapter; they are not
duplicated here.

This route is an exact internal result with explicit starting structure:
it receives \(x_{\rm term}\), its signed reading, \(P,E,\Phi\), the reversible seal,
and the boundary. It does not receive \(\alpha\) or CODATA as input. Metrological
comparison opens only after \eqref{eq:alpha-rev8-salida-corta} has been obtained.

\subsection{Path B: vacancy kernel \texorpdfstring{\(E_{21/22}\)}{E21/22}}

The second route does not invert the charts of \(K\). It starts from the exact density
of vacancy
\[
 \varepsilon_{\rm v}=\log_{10}\!\left(\frac{10}{9}\right)
\]
and of the declared analytic kernel
\begin{equation}
 \mathcal P(x)
 =2\pi x-\frac74x^2+\frac{x^3}{2\pi}
  +\frac{x^4}{20}-\frac{2x^5}{21}-\frac{x^6}{46}.
 \label{eq:alpha-rev8-kernel-e21-22}
\end{equation}
The problem is to solve
\begin{equation}
 \mathcal P(x)=\log_{10}\!\left(\frac{10}{9}\right),
 \qquad 0<x<10^{-2}.
 \label{eq:alpha-rev8-ecuacion-e21-22}
\end{equation}
The theorem of monotonia of the nucleo istablece the uniqueness of the root positive
and the encierro
\[
 0.00729735256928379955
 <x_{21/22}<
 0.00729735256928379957.
\]
Its inverse is
\begin{equation}
 x_{21/22}^{-1}
 =137.0359990840000270200901282598749\ldots.
 \label{eq:alpha-rev8-inversa-e21-22}
\end{equation}
Consequently, the full kernel \(E_{21/22}\), generated by the second
reader of the same enriched state, determines, as a simple and
unique positive root, the same limit section as the dodecaphasic path:
\[
 \boxed{\alpha_A=\alpha_B=\alpha_{\rm HMT}},
 \qquad
 \alpha_B=\operatorname*{arg\,zero}_{x\in I_\alpha}E_{21/22}(x).
\]
The finite rational \(\alpha_{12}\) and the root of the jet \(\mathcal P=P_6\)
are finite projections of the two compatible families; their difference beyond
the printed resolution does not constitute a difference between the two
limit sections. The nested root intervals and the commutative
truncation squares coordinate both prolongations at every depth.

The six coefficients of \(P_6\) are not reconstructed from the root or taken
from a metrological table: they are the first jet of the kernel \(E_{21/22}\) produced
by the internal reader with genealogy \(120/45/11/495\) and prolongation
\(T_{7:9}\). The uniqueness of the root and the internal derivation of the
coefficients are distinct stages of the same already closed path B.

\subsection{Remote prolongation \texorpdfstring{\(\mathsf T_{\alpha,\infty}\)}{T alpha infinity} of the integer path}

\(\mathsf T_{\alpha,\infty}\) extends the recurrence of
\eqref{eq:alpha-rev8-via-entera} to a chain
of length \(N\). To avoid an ambiguity of indices, let
\(\kappa(k)=1+((k-1)\bmod12)\). With a terminal condition at the remote
endpoint one computes
\begin{equation}
\begin{aligned}
 s_k&=P_k+E_k-\Phi_k-K_{\kappa(k)}+c_{k+1},\\
 A_k&=s_k\bmod1000,\\
 c_k&=\left\lfloor\frac{s_k}{1000}\right\rfloor,
 \qquad k=N,N-1,\ldots,1.
\end{aligned}
\label{eq:alpha-rev8-t1}
\end{equation}
The finite case \(N=1000\) uses one thousand triads, that is, three thousand digits. The
carry arriving from the remote end modifies the twelfth
visible block:
\[
\begin{array}{rcl}
\text{short dodecaphasic window, }c_{13}=0
 &:& \ldots|380|663,\\
\text{extension }\mathsf T_{\alpha,\infty}\text{ of one thousand triads}
 &:& \ldots|380|662|999\ldots.
\end{array}
\]
They are not two rival values. They are two boundary semantics of the same
carry law: the first closes at position twelve; the second transports to
it information proceeding from a remote tail.

The reproducible calculation verifies the long word exactly using
\(P,E,\Phi\) as correlated projections of the same joint
APP--TRIT--TPK coinductive state that also publishes the dodecaphasic coordinate
\(\alpha_{\mathrm{HMT}}\). Therefore, that execution certifies carry
propagation: it does not place \(\pi,e,\varphi\) earlier in the genealogy or add
\(\alpha\) from outside. The joint projection
\[
 \mathcal G_9^{\mathrm{joint}}\longrightarrow
 \bigl(\pi,e,\varphi,\mathsf T_{\alpha,\infty}\bigr)
\]
factors through the \cref{sec:alpha-t1-10000-rev9}: the three
Archimedean channels and the dodecaphasic channel open from the same
enriched state; \eqref{eq:alpha-rev8-t1} is then executed and the
independence of the prefix from the five terminal carry
conditions is verified. The separation of projections prevents the long recurrence
from retrospectively selecting its initial data.

\subsection{Causal coordination of dodecaphase closure and the \texorpdfstring{\(E_{21/22}\)}{E21/22} characterization: domains, codomains, and uniqueness}

The classification of this block is fixed as follows:
\begingroup
\raggedright
\begin{description}[style=nextline,leftmargin=3.2cm,labelwidth=2.8cm,labelsep=.4cm]
\item[Generative branch A: dodecaphasic closure.]
It receives the enriched terminal state, the signed chart, the seal, the three
evaluations, and the explicit boundary; it produces the rational
\eqref{eq:alpha-rev8-salida-corta} by means of the integer identity with carry.

\item[analytic-characterization branch: equation \(E_{21/22}\).]
For the declared kernel, monotonicity and the sign change determine a
unique root in \((0,10^{-2})\). The root does not retrospectively determine the
six coefficients of the kernel.

\item[Prolongation \(\mathsf T_{\alpha,\infty}\).]
The long recurrence receives the thousand reference triads, the periodic seal,
and the remote boundary, and transports the carry to the visible window. It does not
constitute the analytic-characterization branch.

\item[\(\mathcal G_9\to\mathsf T_{\alpha,\infty}\) composition.]
Joint rational arithmetic produces ten thousand digits, executes the
carry recurrence, and then verifies prefix stability under the
declared terminal conditions.
\end{description}
\endgroup

The terminal selector, the reversible charts, and the short recurrence are proved in \cref{chap:alpha}; the existence and uniqueness of the second route in \eqref{eq:alpha-rev8-kernel-e21-22}--\eqref{eq:alpha-rev8-inversa-e21-22}; and the long composition in \cref{sec:alpha-t1-10000-rev9}. The computational supplement reproduces the finite enumerations and terminal conditions; the proof rests on the identities and recurrences presented here.

The public causal order is thus closed:
\[
 x_{\rm term}
 \xrightarrow{(\pi_{\rm term},R_{\rm sgn})}
 (B_{\rm term},U_{12}^{\rm sgn})
 \longrightarrow
 \begin{cases}
 \text{path A: dodecaphasic integer closure},\\
 \text{path B: kernel root }E_{21/22},
 \end{cases}
\]
with \(\mathsf T_{\alpha,\infty}\) as remote prolongation of the generative branch and not
as an additional branch. Only
after these outputs is external metrological recognition opened.
